% =========================================================
% Slides Beamer Metropolis 16:9 — Chapter 1 (Basic Statistics)
% Time Series Analysis — Aymen Ben Brik (aymen.benbrik@esprit.tn)
% Esprit School of Business
%
% Migrated from _archives/source-slides/Basic Statistics.tex (Madrid).
% Compile : pdflatex slides.tex (twice).
% =========================================================
\documentclass[aspectratio=169]{beamer}

% =========================================================
% Beamer preamble — Time Series Analysis module
% Theme : Metropolis, aspect ratio 16:9
% Author : Aymen Ben Brik (aymen.benbrik@esprit.tn)
% Esprit School of Business
%
% Include from each chapterX/slides.tex via :
%   % =========================================================
% Beamer preamble — Time Series Analysis module
% Theme : Metropolis, aspect ratio 16:9
% Author : Aymen Ben Brik (aymen.benbrik@esprit.tn)
% Esprit School of Business
%
% Include from each chapterX/slides.tex via :
%   % =========================================================
% Beamer preamble — Time Series Analysis module
% Theme : Metropolis, aspect ratio 16:9
% Author : Aymen Ben Brik (aymen.benbrik@esprit.tn)
% Esprit School of Business
%
% Include from each chapterX/slides.tex via :
%   \input{../preamble/slides-preamble.tex}
% AFTER the line \documentclass[aspectratio=169]{beamer}.
% =========================================================

\usepackage[utf8]{inputenc}
\usepackage[T1]{fontenc}
\usepackage[english]{babel}
\usepackage{lmodern}
\usepackage{amsmath, amssymb, mathtools, bm}
\usepackage{graphicx}
\usepackage{booktabs}
\usepackage{array}
\usepackage{multirow}
\usepackage{colortbl}
\usepackage{xcolor}
\usepackage{tikz}
\usetikzlibrary{positioning, arrows.meta, calc, shapes, fit, backgrounds, matrix, decorations.pathreplacing}
\usepackage{listings}
\usepackage[most]{tcolorbox}

% --- Theme ---
\usetheme{metropolis}
\setbeamertemplate{navigation symbols}{}
\metroset{block=fill}

% --- Esprit colour palette ---
\definecolor{espritBlue}{RGB}{30, 60, 110}
\definecolor{espritAccent}{RGB}{200, 80, 30}
\setbeamercolor{progress bar}{fg=espritAccent}
\setbeamercolor{title separator}{fg=espritAccent}
\setbeamercolor{frametitle}{bg=espritBlue}

% --- Code listings (Python + R, with UTF-8 literate) ---
\definecolor{codebg}{RGB}{248,248,248}
\definecolor{codeframe}{RGB}{220,220,220}
\definecolor{keyword}{RGB}{0,82,155}
\definecolor{comment}{RGB}{88,110,117}
\definecolor{string}{RGB}{133,153,0}

\lstdefinestyle{pythonstyle}{
  language=Python,
  basicstyle=\ttfamily\scriptsize,
  keywordstyle=\color{keyword}\bfseries,
  commentstyle=\color{comment}\itshape,
  stringstyle=\color{string},
  backgroundcolor=\color{codebg},
  frame=single,
  rulecolor=\color{codeframe},
  frameround=tttt,
  breaklines=true,
  showstringspaces=false,
  tabsize=2,
  literate=
    {é}{{\'e}}1 {è}{{\`e}}1 {ê}{{\^e}}1 {à}{{\`a}}1 {â}{{\^a}}1
    {ç}{{\c{c}}}1 {²}{{$^2$}}1 {°}{{${}^\circ$}}1
    {α}{{$\alpha$}}1 {β}{{$\beta$}}1 {γ}{{$\gamma$}}1
    {δ}{{$\delta$}}1 {θ}{{$\theta$}}1 {λ}{{$\lambda$}}1
    {μ}{{$\mu$}}1 {σ}{{$\sigma$}}1 {ε}{{$\varepsilon$}}1
    {ρ}{{$\rho$}}1 {τ}{{$\tau$}}1 {φ}{{$\phi$}}1
    {≤}{{$\leq$}}1 {≥}{{$\geq$}}1 {≈}{{$\approx$}}1
}
\lstdefinestyle{Rstyle}{
  language=R,
  basicstyle=\ttfamily\scriptsize,
  keywordstyle=\color{keyword}\bfseries,
  commentstyle=\color{comment}\itshape,
  stringstyle=\color{string},
  backgroundcolor=\color{codebg},
  frame=single,
  rulecolor=\color{codeframe},
  frameround=tttt,
  breaklines=true,
  showstringspaces=false,
  tabsize=2
}
\lstset{style=pythonstyle}

% --- Common macros ---
\input{../preamble/macros.tex}

% --- tcolorbox boxes for slides ---
\definecolor{boxTh}{RGB}{120, 30, 30}
\definecolor{boxProp}{RGB}{30, 100, 60}
\definecolor{boxDef}{RGB}{30, 60, 110}
\definecolor{boxEx}{RGB}{180, 120, 0}

\newtcolorbox{infobox}[1][]{enhanced, colback=espritBlue!5, colframe=espritBlue,
  fonttitle=\bfseries\small,
  title={\ifstrempty{#1}{}{#1}},
  boxrule=0.4pt, arc=2pt, left=4pt, right=4pt, top=2pt, bottom=2pt}
\newtcolorbox{warnbox}[1][]{enhanced, colback=espritAccent!5, colframe=espritAccent,
  fonttitle=\bfseries\small,
  title={\ifstrempty{#1}{}{#1}},
  boxrule=0.4pt, arc=2pt, left=4pt, right=4pt, top=2pt, bottom=2pt}
\newtcolorbox{thbox}[1][]{enhanced, colback=boxTh!5, colframe=boxTh,
  fonttitle=\bfseries\small,
  title={Theorem\ifstrempty{#1}{}{ -- #1}},
  boxrule=0.5pt, arc=2pt, left=4pt, right=4pt, top=2pt, bottom=2pt}
\newtcolorbox{propbox}[1][]{enhanced, colback=boxProp!5, colframe=boxProp,
  fonttitle=\bfseries\small,
  title={Proposition\ifstrempty{#1}{}{ -- #1}},
  boxrule=0.5pt, arc=2pt, left=4pt, right=4pt, top=2pt, bottom=2pt}
\newtcolorbox{defbox}[1][]{enhanced, colback=boxDef!5, colframe=boxDef,
  fonttitle=\bfseries\small,
  title={Definition\ifstrempty{#1}{}{ -- #1}},
  boxrule=0.5pt, arc=2pt, left=4pt, right=4pt, top=2pt, bottom=2pt}
\newtcolorbox{exbox}[1][]{enhanced, colback=boxEx!5, colframe=boxEx,
  fonttitle=\bfseries\small,
  title={Example\ifstrempty{#1}{}{ -- #1}},
  boxrule=0.5pt, arc=2pt, left=4pt, right=4pt, top=2pt, bottom=2pt}

% AFTER the line \documentclass[aspectratio=169]{beamer}.
% =========================================================

\usepackage[utf8]{inputenc}
\usepackage[T1]{fontenc}
\usepackage[english]{babel}
\usepackage{lmodern}
\usepackage{amsmath, amssymb, mathtools, bm}
\usepackage{graphicx}
\usepackage{booktabs}
\usepackage{array}
\usepackage{multirow}
\usepackage{colortbl}
\usepackage{xcolor}
\usepackage{tikz}
\usetikzlibrary{positioning, arrows.meta, calc, shapes, fit, backgrounds, matrix, decorations.pathreplacing}
\usepackage{listings}
\usepackage[most]{tcolorbox}

% --- Theme ---
\usetheme{metropolis}
\setbeamertemplate{navigation symbols}{}
\metroset{block=fill}

% --- Esprit colour palette ---
\definecolor{espritBlue}{RGB}{30, 60, 110}
\definecolor{espritAccent}{RGB}{200, 80, 30}
\setbeamercolor{progress bar}{fg=espritAccent}
\setbeamercolor{title separator}{fg=espritAccent}
\setbeamercolor{frametitle}{bg=espritBlue}

% --- Code listings (Python + R, with UTF-8 literate) ---
\definecolor{codebg}{RGB}{248,248,248}
\definecolor{codeframe}{RGB}{220,220,220}
\definecolor{keyword}{RGB}{0,82,155}
\definecolor{comment}{RGB}{88,110,117}
\definecolor{string}{RGB}{133,153,0}

\lstdefinestyle{pythonstyle}{
  language=Python,
  basicstyle=\ttfamily\scriptsize,
  keywordstyle=\color{keyword}\bfseries,
  commentstyle=\color{comment}\itshape,
  stringstyle=\color{string},
  backgroundcolor=\color{codebg},
  frame=single,
  rulecolor=\color{codeframe},
  frameround=tttt,
  breaklines=true,
  showstringspaces=false,
  tabsize=2,
  literate=
    {é}{{\'e}}1 {è}{{\`e}}1 {ê}{{\^e}}1 {à}{{\`a}}1 {â}{{\^a}}1
    {ç}{{\c{c}}}1 {²}{{$^2$}}1 {°}{{${}^\circ$}}1
    {α}{{$\alpha$}}1 {β}{{$\beta$}}1 {γ}{{$\gamma$}}1
    {δ}{{$\delta$}}1 {θ}{{$\theta$}}1 {λ}{{$\lambda$}}1
    {μ}{{$\mu$}}1 {σ}{{$\sigma$}}1 {ε}{{$\varepsilon$}}1
    {ρ}{{$\rho$}}1 {τ}{{$\tau$}}1 {φ}{{$\phi$}}1
    {≤}{{$\leq$}}1 {≥}{{$\geq$}}1 {≈}{{$\approx$}}1
}
\lstdefinestyle{Rstyle}{
  language=R,
  basicstyle=\ttfamily\scriptsize,
  keywordstyle=\color{keyword}\bfseries,
  commentstyle=\color{comment}\itshape,
  stringstyle=\color{string},
  backgroundcolor=\color{codebg},
  frame=single,
  rulecolor=\color{codeframe},
  frameround=tttt,
  breaklines=true,
  showstringspaces=false,
  tabsize=2
}
\lstset{style=pythonstyle}

% --- Common macros ---
% =========================================================
% Common math macros — Time Series Analysis module
% Author : Aymen Ben Brik (aymen.benbrik@esprit.tn)
% Esprit School of Business
% =========================================================

% --- Standard sets ---
\providecommand{\R}{\mathbb{R}}
\providecommand{\N}{\mathbb{N}}
\providecommand{\Z}{\mathbb{Z}}
\providecommand{\Q}{\mathbb{Q}}
\providecommand{\C}{\mathbb{C}}

% --- Probability / statistics ---
\providecommand{\E}{\mathbb{E}}
\providecommand{\Prob}{\mathbb{P}}
\providecommand{\Var}{\mathrm{Var}}
\providecommand{\Cov}{\mathrm{Cov}}
\providecommand{\Corr}{\mathrm{Corr}}
\providecommand{\indep}{\perp\!\!\!\perp}
\providecommand{\iid}{\overset{\text{i.i.d.}}{\sim}}
\providecommand{\given}{\,\big\vert\,}

% --- Estimators / hat ---
\providecommand{\hatm}{\widehat{m}}
\providecommand{\hatsig}{\widehat{\sigma}}
\providecommand{\hatrho}{\widehat{\rho}}
\providecommand{\hatphi}{\widehat{\phi}}
\providecommand{\hattheta}{\widehat{\theta}}

% --- Time series notation ---
\providecommand{\Lop}{L}                      % lag operator
\providecommand{\Diff}{\Delta}                % difference operator
\providecommand{\AR}[1]{\mathrm{AR}(#1)}
\providecommand{\MAm}[1]{\mathrm{MA}(#1)}
\providecommand{\ARMA}[2]{\mathrm{ARMA}(#1,#2)}
\providecommand{\ARIMA}[3]{\mathrm{ARIMA}(#1,#2,#3)}
\providecommand{\SARIMA}[6]{\mathrm{SARIMA}(#1,#2,#3)(#4,#5,#6)}
\providecommand{\ARCH}[1]{\mathrm{ARCH}(#1)}
\providecommand{\GARCH}[2]{\mathrm{GARCH}(#1,#2)}
\providecommand{\WN}{\mathrm{WN}}              % white noise
\providecommand{\MSE}{\mathrm{MSE}}
\providecommand{\MAE}{\mathrm{MAE}}
\providecommand{\RMSE}{\mathrm{RMSE}}
\providecommand{\AIC}{\mathrm{AIC}}
\providecommand{\BIC}{\mathrm{BIC}}
\providecommand{\ACF}{\mathrm{ACF}}
\providecommand{\PACF}{\mathrm{PACF}}

% --- Linear algebra ---
\providecommand{\transp}[1]{#1^{\!\top}}
\providecommand{\norm}[1]{\left\lVert #1 \right\rVert}
\providecommand{\abs}[1]{\left\lvert #1 \right\rvert}
\providecommand{\inner}[2]{\langle #1, #2 \rangle}
\providecommand{\diag}{\mathrm{diag}}
\providecommand{\tr}{\mathrm{tr}}

% --- Bold vectors / matrices ---
\providecommand{\vx}{\bm{x}}
\providecommand{\vy}{\bm{y}}
\providecommand{\vz}{\bm{z}}
\providecommand{\vh}{\bm{h}}
\providecommand{\vu}{\bm{u}}
\providecommand{\vv}{\bm{v}}
\providecommand{\vw}{\bm{w}}
\providecommand{\vb}{\bm{b}}
\providecommand{\vW}{\bm{W}}
\providecommand{\vSigma}{\bm{\Sigma}}
\providecommand{\veps}{\bm{\varepsilon}}

% --- Author identity (reused on titlepages) ---
\providecommand{\auteurNom}{Aymen Ben Brik}
\providecommand{\auteurEmail}{aymen.benbrik@esprit.tn}
\providecommand{\auteurInstitution}{Esprit School of Business}
\providecommand{\auteurRole}{Module head}
\providecommand{\moduleNom}{Time Series Analysis}
\providecommand{\anneeUniversitaire}{2025--2026}


% --- tcolorbox boxes for slides ---
\definecolor{boxTh}{RGB}{120, 30, 30}
\definecolor{boxProp}{RGB}{30, 100, 60}
\definecolor{boxDef}{RGB}{30, 60, 110}
\definecolor{boxEx}{RGB}{180, 120, 0}

\newtcolorbox{infobox}[1][]{enhanced, colback=espritBlue!5, colframe=espritBlue,
  fonttitle=\bfseries\small,
  title={\ifstrempty{#1}{}{#1}},
  boxrule=0.4pt, arc=2pt, left=4pt, right=4pt, top=2pt, bottom=2pt}
\newtcolorbox{warnbox}[1][]{enhanced, colback=espritAccent!5, colframe=espritAccent,
  fonttitle=\bfseries\small,
  title={\ifstrempty{#1}{}{#1}},
  boxrule=0.4pt, arc=2pt, left=4pt, right=4pt, top=2pt, bottom=2pt}
\newtcolorbox{thbox}[1][]{enhanced, colback=boxTh!5, colframe=boxTh,
  fonttitle=\bfseries\small,
  title={Theorem\ifstrempty{#1}{}{ -- #1}},
  boxrule=0.5pt, arc=2pt, left=4pt, right=4pt, top=2pt, bottom=2pt}
\newtcolorbox{propbox}[1][]{enhanced, colback=boxProp!5, colframe=boxProp,
  fonttitle=\bfseries\small,
  title={Proposition\ifstrempty{#1}{}{ -- #1}},
  boxrule=0.5pt, arc=2pt, left=4pt, right=4pt, top=2pt, bottom=2pt}
\newtcolorbox{defbox}[1][]{enhanced, colback=boxDef!5, colframe=boxDef,
  fonttitle=\bfseries\small,
  title={Definition\ifstrempty{#1}{}{ -- #1}},
  boxrule=0.5pt, arc=2pt, left=4pt, right=4pt, top=2pt, bottom=2pt}
\newtcolorbox{exbox}[1][]{enhanced, colback=boxEx!5, colframe=boxEx,
  fonttitle=\bfseries\small,
  title={Example\ifstrempty{#1}{}{ -- #1}},
  boxrule=0.5pt, arc=2pt, left=4pt, right=4pt, top=2pt, bottom=2pt}

% AFTER the line \documentclass[aspectratio=169]{beamer}.
% =========================================================

\usepackage[utf8]{inputenc}
\usepackage[T1]{fontenc}
\usepackage[english]{babel}
\usepackage{lmodern}
\usepackage{amsmath, amssymb, mathtools, bm}
\usepackage{graphicx}
\usepackage{booktabs}
\usepackage{array}
\usepackage{multirow}
\usepackage{colortbl}
\usepackage{xcolor}
\usepackage{tikz}
\usetikzlibrary{positioning, arrows.meta, calc, shapes, fit, backgrounds, matrix, decorations.pathreplacing}
\usepackage{listings}
\usepackage[most]{tcolorbox}

% --- Theme ---
\usetheme{metropolis}
\setbeamertemplate{navigation symbols}{}
\metroset{block=fill}

% --- Esprit colour palette ---
\definecolor{espritBlue}{RGB}{30, 60, 110}
\definecolor{espritAccent}{RGB}{200, 80, 30}
\setbeamercolor{progress bar}{fg=espritAccent}
\setbeamercolor{title separator}{fg=espritAccent}
\setbeamercolor{frametitle}{bg=espritBlue}

% --- Code listings (Python + R, with UTF-8 literate) ---
\definecolor{codebg}{RGB}{248,248,248}
\definecolor{codeframe}{RGB}{220,220,220}
\definecolor{keyword}{RGB}{0,82,155}
\definecolor{comment}{RGB}{88,110,117}
\definecolor{string}{RGB}{133,153,0}

\lstdefinestyle{pythonstyle}{
  language=Python,
  basicstyle=\ttfamily\scriptsize,
  keywordstyle=\color{keyword}\bfseries,
  commentstyle=\color{comment}\itshape,
  stringstyle=\color{string},
  backgroundcolor=\color{codebg},
  frame=single,
  rulecolor=\color{codeframe},
  frameround=tttt,
  breaklines=true,
  showstringspaces=false,
  tabsize=2,
  literate=
    {é}{{\'e}}1 {è}{{\`e}}1 {ê}{{\^e}}1 {à}{{\`a}}1 {â}{{\^a}}1
    {ç}{{\c{c}}}1 {²}{{$^2$}}1 {°}{{${}^\circ$}}1
    {α}{{$\alpha$}}1 {β}{{$\beta$}}1 {γ}{{$\gamma$}}1
    {δ}{{$\delta$}}1 {θ}{{$\theta$}}1 {λ}{{$\lambda$}}1
    {μ}{{$\mu$}}1 {σ}{{$\sigma$}}1 {ε}{{$\varepsilon$}}1
    {ρ}{{$\rho$}}1 {τ}{{$\tau$}}1 {φ}{{$\phi$}}1
    {≤}{{$\leq$}}1 {≥}{{$\geq$}}1 {≈}{{$\approx$}}1
}
\lstdefinestyle{Rstyle}{
  language=R,
  basicstyle=\ttfamily\scriptsize,
  keywordstyle=\color{keyword}\bfseries,
  commentstyle=\color{comment}\itshape,
  stringstyle=\color{string},
  backgroundcolor=\color{codebg},
  frame=single,
  rulecolor=\color{codeframe},
  frameround=tttt,
  breaklines=true,
  showstringspaces=false,
  tabsize=2
}
\lstset{style=pythonstyle}

% --- Common macros ---
% =========================================================
% Common math macros — Time Series Analysis module
% Author : Aymen Ben Brik (aymen.benbrik@esprit.tn)
% Esprit School of Business
% =========================================================

% --- Standard sets ---
\providecommand{\R}{\mathbb{R}}
\providecommand{\N}{\mathbb{N}}
\providecommand{\Z}{\mathbb{Z}}
\providecommand{\Q}{\mathbb{Q}}
\providecommand{\C}{\mathbb{C}}

% --- Probability / statistics ---
\providecommand{\E}{\mathbb{E}}
\providecommand{\Prob}{\mathbb{P}}
\providecommand{\Var}{\mathrm{Var}}
\providecommand{\Cov}{\mathrm{Cov}}
\providecommand{\Corr}{\mathrm{Corr}}
\providecommand{\indep}{\perp\!\!\!\perp}
\providecommand{\iid}{\overset{\text{i.i.d.}}{\sim}}
\providecommand{\given}{\,\big\vert\,}

% --- Estimators / hat ---
\providecommand{\hatm}{\widehat{m}}
\providecommand{\hatsig}{\widehat{\sigma}}
\providecommand{\hatrho}{\widehat{\rho}}
\providecommand{\hatphi}{\widehat{\phi}}
\providecommand{\hattheta}{\widehat{\theta}}

% --- Time series notation ---
\providecommand{\Lop}{L}                      % lag operator
\providecommand{\Diff}{\Delta}                % difference operator
\providecommand{\AR}[1]{\mathrm{AR}(#1)}
\providecommand{\MAm}[1]{\mathrm{MA}(#1)}
\providecommand{\ARMA}[2]{\mathrm{ARMA}(#1,#2)}
\providecommand{\ARIMA}[3]{\mathrm{ARIMA}(#1,#2,#3)}
\providecommand{\SARIMA}[6]{\mathrm{SARIMA}(#1,#2,#3)(#4,#5,#6)}
\providecommand{\ARCH}[1]{\mathrm{ARCH}(#1)}
\providecommand{\GARCH}[2]{\mathrm{GARCH}(#1,#2)}
\providecommand{\WN}{\mathrm{WN}}              % white noise
\providecommand{\MSE}{\mathrm{MSE}}
\providecommand{\MAE}{\mathrm{MAE}}
\providecommand{\RMSE}{\mathrm{RMSE}}
\providecommand{\AIC}{\mathrm{AIC}}
\providecommand{\BIC}{\mathrm{BIC}}
\providecommand{\ACF}{\mathrm{ACF}}
\providecommand{\PACF}{\mathrm{PACF}}

% --- Linear algebra ---
\providecommand{\transp}[1]{#1^{\!\top}}
\providecommand{\norm}[1]{\left\lVert #1 \right\rVert}
\providecommand{\abs}[1]{\left\lvert #1 \right\rvert}
\providecommand{\inner}[2]{\langle #1, #2 \rangle}
\providecommand{\diag}{\mathrm{diag}}
\providecommand{\tr}{\mathrm{tr}}

% --- Bold vectors / matrices ---
\providecommand{\vx}{\bm{x}}
\providecommand{\vy}{\bm{y}}
\providecommand{\vz}{\bm{z}}
\providecommand{\vh}{\bm{h}}
\providecommand{\vu}{\bm{u}}
\providecommand{\vv}{\bm{v}}
\providecommand{\vw}{\bm{w}}
\providecommand{\vb}{\bm{b}}
\providecommand{\vW}{\bm{W}}
\providecommand{\vSigma}{\bm{\Sigma}}
\providecommand{\veps}{\bm{\varepsilon}}

% --- Author identity (reused on titlepages) ---
\providecommand{\auteurNom}{Aymen Ben Brik}
\providecommand{\auteurEmail}{aymen.benbrik@esprit.tn}
\providecommand{\auteurInstitution}{Esprit School of Business}
\providecommand{\auteurRole}{Module head}
\providecommand{\moduleNom}{Time Series Analysis}
\providecommand{\anneeUniversitaire}{2025--2026}


% --- tcolorbox boxes for slides ---
\definecolor{boxTh}{RGB}{120, 30, 30}
\definecolor{boxProp}{RGB}{30, 100, 60}
\definecolor{boxDef}{RGB}{30, 60, 110}
\definecolor{boxEx}{RGB}{180, 120, 0}

\newtcolorbox{infobox}[1][]{enhanced, colback=espritBlue!5, colframe=espritBlue,
  fonttitle=\bfseries\small,
  title={\ifstrempty{#1}{}{#1}},
  boxrule=0.4pt, arc=2pt, left=4pt, right=4pt, top=2pt, bottom=2pt}
\newtcolorbox{warnbox}[1][]{enhanced, colback=espritAccent!5, colframe=espritAccent,
  fonttitle=\bfseries\small,
  title={\ifstrempty{#1}{}{#1}},
  boxrule=0.4pt, arc=2pt, left=4pt, right=4pt, top=2pt, bottom=2pt}
\newtcolorbox{thbox}[1][]{enhanced, colback=boxTh!5, colframe=boxTh,
  fonttitle=\bfseries\small,
  title={Theorem\ifstrempty{#1}{}{ -- #1}},
  boxrule=0.5pt, arc=2pt, left=4pt, right=4pt, top=2pt, bottom=2pt}
\newtcolorbox{propbox}[1][]{enhanced, colback=boxProp!5, colframe=boxProp,
  fonttitle=\bfseries\small,
  title={Proposition\ifstrempty{#1}{}{ -- #1}},
  boxrule=0.5pt, arc=2pt, left=4pt, right=4pt, top=2pt, bottom=2pt}
\newtcolorbox{defbox}[1][]{enhanced, colback=boxDef!5, colframe=boxDef,
  fonttitle=\bfseries\small,
  title={Definition\ifstrempty{#1}{}{ -- #1}},
  boxrule=0.5pt, arc=2pt, left=4pt, right=4pt, top=2pt, bottom=2pt}
\newtcolorbox{exbox}[1][]{enhanced, colback=boxEx!5, colframe=boxEx,
  fonttitle=\bfseries\small,
  title={Example\ifstrempty{#1}{}{ -- #1}},
  boxrule=0.5pt, arc=2pt, left=4pt, right=4pt, top=2pt, bottom=2pt}


\title[\moduleNom\ -- Ch.1]{Basic Statistical Concepts}
\subtitle{Moments $\cdot$ Estimators $\cdot$ MoM \& MLE $\cdot$ Inference}
\author{\auteurNom \\ \footnotesize\auteurEmail}
\institute{\auteurInstitution}
\date{\anneeUniversitaire}

\begin{document}

% ============== Title & outline ==============
\begin{frame}\titlepage\end{frame}

\begin{frame}{Learning objectives}
  By the end of this chapter, you will be able to:
  \begin{itemize}
    \item interpret the \emph{moments} of a distribution
          (mean, variance, skewness, kurtosis);
    \item compute classical \emph{point estimators} (sample mean,
          variance, ACF) and discuss their properties (bias, consistency);
    \item apply two estimation methods: \emph{Method of Moments} and
          \emph{Maximum Likelihood};
    \item describe the \emph{sampling distribution} of standard
          estimators under normality;
    \item run a basic \emph{statistical inference} procedure
          (confidence interval, hypothesis test, $p$-value).
  \end{itemize}
\end{frame}

\begin{frame}{Chapter outline}
  \begin{enumerate}
    \item Moments of a distribution
    \item Classical estimators and sampling distributions
    \item Estimation methods: MoM and MLE
    \item Multivariate data
    \item Statistical inference (CI, hypothesis testing)
  \end{enumerate}
\end{frame}

% =====================================================================
\section{Moments of a distribution}

\begin{frame}{1.~Raw and central moments}
  \begin{defbox}[Moments]
    Let $X$ be a random variable with density $f$ and mean
    $\mu = \E[X]$.
    \begin{itemize}
      \item \textbf{Raw $\ell$-th moment:}
            $m'_\ell = \E[X^\ell] = \int x^\ell f(x)\, dx$.
      \item \textbf{Central $\ell$-th moment:}
            $m_\ell = \E[(X - \mu)^\ell]$.
    \end{itemize}
  \end{defbox}
  \medskip
  \textbf{Standard quantities:}
  \[
    \mu = m'_1, \qquad
    \sigma^2 = m_2 = \Var[X], \qquad
    S(X) = \frac{m_3}{\sigma^3}, \qquad
    K(X) = \frac{m_4}{\sigma^4}.
  \]
\end{frame}

\begin{frame}{1.~Numerical example: $X \sim \mathcal{N}(2, 4)$}
  Theoretical moments:
  \[
    \mu = 2, \quad \sigma^2 = 4, \quad S(X) = 0, \quad K(X) = 3.
  \]
  \medskip
  Empirical estimates from $n = 10\,000$ Monte-Carlo draws (seed 42):
  \[
    \widehat{\mu} \approx 2.005, \quad
    \widehat{\sigma}^2 \approx 3.987, \quad
    \widehat{S} \approx -0.013, \quad
    \widehat{K} \approx 2.996.
  \]
  \medskip
  \begin{infobox}
    \emph{Skewness} measures asymmetry; \emph{kurtosis} measures
    tail-heaviness. A normal distribution has $S = 0$ and $K = 3$
    (excess kurtosis $K - 3 = 0$).
  \end{infobox}
\end{frame}

% =====================================================================
\section{Classical estimators}

\begin{frame}{2.~Sample mean and sample variance}
  Given i.i.d.\ data $X_1, \dots, X_n$:
  \[
    \bar{X}_n = \frac{1}{n}\sum_{i=1}^{n} X_i,
    \qquad
    S_n^2 = \frac{1}{n-1}\sum_{i=1}^{n}(X_i - \bar{X}_n)^2.
  \]
  \medskip
  \begin{propbox}[Properties]
    Under i.i.d.\ assumption with finite second moments:
    \begin{itemize}
      \item \textbf{Unbiased:} $\E[\bar{X}_n] = \mu$ and $\E[S_n^2] = \sigma^2$.
      \item \textbf{Consistent:} $\bar{X}_n \xrightarrow{p} \mu$
            and $S_n^2 \xrightarrow{p} \sigma^2$ (Law of Large Numbers).
    \end{itemize}
  \end{propbox}
\end{frame}

\begin{frame}{2.~Sampling distribution under normality}
  If $X_1, \dots, X_n \iid \mathcal{N}(\mu, \sigma^2)$:
  \begin{thbox}
    \[
      \bar{X}_n \sim \mathcal{N}\!\Bigl(\mu,\, \tfrac{\sigma^2}{n}\Bigr),
      \qquad
      \frac{(n-1)\, S_n^2}{\sigma^2} \sim \chi^2_{n-1},
      \qquad
      \frac{\bar{X}_n - \mu}{S_n / \sqrt{n}} \sim t_{n-1}.
    \]
  \end{thbox}
  \medskip
  \begin{itemize}
    \item $\bar{X}_n$ and $S_n^2$ are \emph{independent} under normality
          (Cochran's theorem).
    \item These distributions are the basis of $z$-tests, $t$-tests
          and $\chi^2$-tests.
  \end{itemize}
\end{frame}

\begin{frame}{2.~Central Limit Theorem}
  \begin{thbox}[CLT]
    For i.i.d.\ $X_1, \dots, X_n$ with $\E[X_i] = \mu$ and
    $\Var[X_i] = \sigma^2 < \infty$:
    \[
      \frac{\bar{X}_n - \mu}{\sigma / \sqrt{n}}
      \;\xrightarrow{d}\; \mathcal{N}(0, 1).
    \]
  \end{thbox}
  \medskip
  \begin{infobox}
    Even when $X_i$ is \emph{not} normal, the sample mean is
    approximately normal for large $n$. Rule of thumb: $n \geq 30$.
  \end{infobox}
\end{frame}

% =====================================================================
\section{Estimation methods}

\begin{frame}{3.~Method of Moments (MoM)}
  \begin{defbox}[MoM]
    To estimate $\theta \in \R^k$, solve the system
    \[
      \E_\theta[X^p] \;=\; \tfrac{1}{n}\sum_{i=1}^{n} X_i^p,
      \qquad p = 1, \dots, k.
    \]
  \end{defbox}
  \medskip
  \textbf{Example: $\mathcal{N}(\mu, \sigma^2)$.}
  \[
    \mu = \E[X] \;\Rightarrow\; \widehat{\mu}_{\text{MoM}} = \bar{X}_n,
  \]
  \[
    \sigma^2 = \E[X^2] - \mu^2 \;\Rightarrow\;
    \widehat{\sigma}^2_{\text{MoM}} = \tfrac{1}{n}\sum(X_i - \bar{X}_n)^2.
  \]
  Note: this differs from $S_n^2$ by a factor $n / (n-1)$.
\end{frame}

\begin{frame}{3.~Maximum Likelihood (MLE)}
  \begin{defbox}[Likelihood and MLE]
    Likelihood:
    $L(\theta) = \prod_{i=1}^{n} f(x_i; \theta)$ for i.i.d.\ data.
    Log-likelihood: $\ell(\theta) = \log L(\theta)$.
    The MLE is
    \[
      \widehat{\theta}_{\text{MLE}} = \arg\max_\theta\, \ell(\theta).
    \]
  \end{defbox}
  \medskip
  \textbf{Example: $\mathcal{N}(\mu, \sigma^2)$.}
  Solving $\partial \ell / \partial \mu = 0$ and
  $\partial \ell / \partial \sigma^2 = 0$:
  \[
    \widehat{\mu}_{\text{MLE}} = \bar{X}_n,
    \qquad
    \widehat{\sigma}^2_{\text{MLE}} = \tfrac{1}{n}\sum(X_i - \bar{X}_n)^2.
  \]
  Same as MoM for the Gaussian. Both are biased for $\sigma^2$.
\end{frame}

\begin{frame}{3.~Asymptotic normality of the MLE}
  \begin{thbox}
    Under regularity conditions (identifiability, smoothness, finite
    Fisher information $I(\theta_0)$):
    \[
      \sqrt{n}\,\bigl(\widehat{\theta}_n - \theta_0\bigr)
      \;\xrightarrow{d}\;
      \mathcal{N}\!\bigl(0,\, I(\theta_0)^{-1}\bigr).
    \]
  \end{thbox}
  \medskip
  \begin{itemize}
    \item MLE is \textbf{asymptotically efficient} (reaches the
          Cramér–Rao bound).
    \item MoM is simpler but generally less efficient.
    \item Both are consistent.
  \end{itemize}
\end{frame}

% =====================================================================
\section{Multivariate data}

\begin{frame}{4.~Joint, marginal, conditional}
  Joint density: $f(x, y)$.
  \[
    f_X(x) = \int f(x, y)\, dy, \quad
    f_{Y \mid X}(y \mid x) = \frac{f(x, y)}{f_X(x)},
  \]
  \[
    f(x, y) = f_{Y \mid X}(y \mid x)\, f_X(x).
  \]
  \medskip
  \textbf{Independence:} $f(x, y) = f_X(x)\, f_Y(y)$.
  \[
    \E[XY] = \E[X]\E[Y], \qquad
    \Cov(X, Y) = 0.
  \]
\end{frame}

\begin{frame}{4.~Sample covariance and correlation}
  Given paired data $(X_i, Y_i)_{i=1}^{n}$:
  \[
    \widehat{\Cov}(X, Y) = \frac{1}{n-1}\sum_{i=1}^{n}
      (X_i - \bar{X})(Y_i - \bar{Y}),
  \]
  \[
    \widehat{\Corr}(X, Y) =
      \frac{\widehat{\Cov}(X, Y)}{\widehat{\sigma}_X\, \widehat{\sigma}_Y}
      \;\in\; [-1, 1].
  \]
  \medskip
  \begin{warnbox}
    Correlation captures \emph{linear} dependence only. Two variables
    can have $\Corr(X, Y) = 0$ and still be strongly dependent
    (e.g.\ $Y = X^2$ with $X$ symmetric).
  \end{warnbox}
\end{frame}

% =====================================================================
\section{Statistical inference}

\begin{frame}{5.~Confidence interval for the mean}
  Setting: $X_1, \dots, X_n \iid \mathcal{N}(\mu, \sigma^2)$.
  \medskip
  \textbf{$\sigma$ known} ($z$-interval):
  \[
    \mathrm{CI}_{1-\alpha}(\mu)
    = \Bigl[\,\bar{X}_n - z_{1-\alpha/2}\, \tfrac{\sigma}{\sqrt{n}},\;
            \bar{X}_n + z_{1-\alpha/2}\, \tfrac{\sigma}{\sqrt{n}}\Bigr].
  \]
  \textbf{$\sigma$ unknown} ($t$-interval):
  \[
    \mathrm{CI}_{1-\alpha}(\mu)
    = \Bigl[\,\bar{X}_n - t_{n-1,1-\alpha/2}\, \tfrac{S_n}{\sqrt{n}},\;
            \bar{X}_n + t_{n-1,1-\alpha/2}\, \tfrac{S_n}{\sqrt{n}}\Bigr].
  \]
  \medskip
  \textbf{Reading:} probability $1 - \alpha$ that the true $\mu$ falls
  in the interval (over repeated samples).
\end{frame}

\begin{frame}{5.~Hypothesis testing}
  \begin{defbox}[Two-sided test on the mean]
    $H_0: \mu = \mu_0$ vs.\ $H_1: \mu \neq \mu_0$. Test statistic
    $T = (\bar{X}_n - \mu_0) / (S_n / \sqrt{n}) \sim t_{n-1}$ under $H_0$.
    Reject $H_0$ at level $\alpha$ if $|T| > t_{n-1, 1-\alpha/2}$.
  \end{defbox}
  \medskip
  \begin{itemize}
    \item \textbf{$p$-value:} probability under $H_0$ of observing
          something at least as extreme as the data.
    \item \textbf{Type I error} ($\alpha$): rejecting $H_0$ when true.
    \item \textbf{Type II error} ($\beta$): not rejecting $H_0$ when false.
    \item \textbf{Power:} $1 - \beta$ — sensitivity to detect $H_1$.
  \end{itemize}
\end{frame}

\begin{frame}{5.~Numerical example}
  We sample $n = 50$ Atlanta monthly temperatures (in $^\circ$F).
  Sample statistics: $\bar{X}_{50} = 61.2$, $S_{50} = 16.8$.
  Test $H_0: \mu = 60$ vs.\ $H_1: \mu \neq 60$ at $\alpha = 0.05$.
  \medskip
  \[
    T = \frac{61.2 - 60}{16.8 / \sqrt{50}} \approx 0.505,
    \qquad
    t_{49, 0.975} \approx 2.010.
  \]
  Since $|T| = 0.505 < 2.010$, \textbf{do not reject} $H_0$.
  $p$-value $\approx 0.616$.
  \medskip
  \begin{infobox}
    The 95\% CI for $\mu$ is
    $[61.2 - 4.78,\; 61.2 + 4.78] = [56.4,\, 66.0]$,
    which contains $60$ — consistent with the test conclusion.
  \end{infobox}
\end{frame}

% =====================================================================
\section{Wrap-up}

\begin{frame}{Key takeaways}
  \begin{itemize}
    \item The \emph{moments} characterise a distribution
          (mean, variance, skewness, kurtosis).
    \item \emph{Classical estimators} ($\bar{X}, S^2$) are unbiased
          and consistent under i.i.d.\ assumptions.
    \item \emph{MoM} and \emph{MLE} coincide for the Gaussian; in
          general MLE is asymptotically efficient, MoM is simpler.
    \item Under normality, $\bar{X} \sim \mathcal{N}(\mu, \sigma^2/n)$
          and $(n-1)S^2 / \sigma^2 \sim \chi^2_{n-1}$.
    \item \emph{CLT} extends asymptotic normality to non-Gaussian data.
    \item Confidence intervals and hypothesis tests are
          built on the sampling distribution of the estimator.
  \end{itemize}
\end{frame}

\begin{frame}{Up next}
  \textbf{Chapter 2.} Time-series \emph{decomposition}: trend and
  seasonality.
  \medskip
  \begin{itemize}
    \item Classical additive and multiplicative decomposition.
    \item Trend estimation by moving averages and parametric models.
    \item Seasonality estimation.
    \item Residual analysis.
  \end{itemize}
  \medskip
  \begin{center}\Large\bfseries Thank you.\end{center}
  \begin{center}\footnotesize\auteurNom\ -- \auteurEmail\ -- \auteurInstitution\end{center}
\end{frame}

\end{document}
